% === Section 12: Limitations ===
\section{Limitations and Future Work}\label{sec:limitations}

\subsection{Limitations}

\paragraph{Two-domain scope.}
This paper validates ORIUS on two CPS domains.  While the kernel is designed to
be universal, the empirical evidence covers only battery energy storage and
autonomous vehicles.  Four additional candidate domains (industrial monitoring,
healthcare, maritime navigation, aerospace) motivated the architecture but have
\textbf{no locked pipeline artifacts} in this release and are not the basis of any
empirical claim.  Validation on additional domains is required to strengthen the
universality claim; until then, T11 (Typed Structural Transfer) is proven
formally but confirmed empirically only for Battery $\leftrightarrow$ AV.

\paragraph{AV TSVR floor.}
The AV domain achieves a TSVR of 0.628, which, while improved over the baseline
(0.660), is still substantial.  The residual violations arise primarily from the 4\,s
prediction horizon, where complex multi-agent interactions produce true states
outside even the widened conformal interval.  A more expressive predictor (\eg~a
transformer-based trajectory model) could lower the base TSVR.  The AV result should
be read as a \emph{bounded proof-of-mechanism}, not as a claim of zero-violation
closure.

\paragraph{Conformal coverage assumptions.}
The finite-sample coverage guarantee of split conformal prediction holds
\emph{marginally} (averaged over the calibration set) but not \emph{conditionally}
(for every individual input).  The shift-aware mechanism provides robustness against
subgroup-level violations but does not eliminate them entirely.  Per-input guarantees
would require conditional conformal methods not yet integrated.

\paragraph{Shift guarantees are bounded.}
The shift-aware adaptive mechanism (Section~\ref{sec:shift}) compensates for
distribution drift within the observed rate during validation (A6).  It has
\textbf{not} been evaluated under intentionally adversarial telemetry attacks,
sudden-onset covariate shift, or concept drift that violates the bounded-rate
assumption.  Guarantees degrade gracefully but may fail silently under pathological
shift patterns.

\paragraph{Fallback action conservatism.}
Assumption A4 requires that the fallback action is safe for all states.  In many
domains, designing a universally safe fallback is straightforward (\eg~zero
dispatch for battery, maximum braking for AV).  In domains with coupled dynamics,
identifying a universally safe fallback may be non-trivial.

\paragraph{Single-sensor pipeline.}
The current detect stage focuses on single-channel telemetry faults.  Cross-sensor
consistency checks (\eg~lidar--camera fusion disagreement) are not yet implemented
in the adapter interface.

\paragraph{Computational cost.}
The AV runtime requires $\sim$55~minutes for 9{,}348~steps on a single CPU core.
For real-time deployment at 10~Hz, the per-step budget is 100~ms, whereas the
current implementation averages $\sim$350~ms.  Optimisation (C++ inference,
batched prediction, GPU acceleration) would be needed for real-time use.

\paragraph{Cross-domain transfer scope.}
The adapter interface enables structural reuse across domains, but functional
equivalence (whether the same hyperparameters, calibration windows, and inflation
schedules are optimal) has not been demonstrated beyond Battery and AV.  Adapting
to a new domain still requires implementing the six typed adapter methods and
conducting a domain-specific calibration run.

\subsection{Future Work}

\begin{enumerate}[nosep]
  \item \textbf{Additional domains}: Extend validation to industrial process
        control, healthcare monitoring, and autonomous maritime systems.
  \item \textbf{Stronger predictors}: Replace LightGBM with transformer-based
        trajectory predictors to lower the base TSVR in the AV domain.
  \item \textbf{Conditional coverage}: Integrate conditional conformal prediction
        methods (\eg~locally-adaptive intervals) for per-input guarantees.
  \item \textbf{Multi-sensor fusion}: Extend the detect stage to perform
        cross-sensor consistency checks using the adapter interface.
  \item \textbf{Real-time deployment}: Optimise the runtime for sub-100~ms latency
        via compiled inference and hardware acceleration.
  \item \textbf{Formal verification}: Apply model-checking or SMT-based verification
        to the kernel contracts for certified correctness beyond testing.
  \item \textbf{Adversarial robustness}: Evaluate ORIUS under intentionally
        adversarial telemetry attacks beyond the fault-injection families tested here.
\end{enumerate}
